\subsection{位}

若 K 是域，回忆 \(\tilde{K}\) 表示 K 与一个记作 \(\infty\) 的元素之和（的集）（《代数》第二章 § 9 第 9 节）；K 的合成律按如下方式（出处同上）扩张到 \(\tilde{K}\)

\begin{enumerate}
\def\labelenumi{(\arabic{enumi})}
\setcounter{enumi}{1}
\tightlist
\item
\end{enumerate}

\[
a + \infty = \infty \quad \text { for } \quad a \in \mathbf {K}, \quad a \neq \infty ,\tag{2}
\]

\[
\infty . a = a. \infty = \infty \quad \text {   for   } \quad a \in \tilde {\mathbf {K}}, \quad a \neq 0.\tag{3}
\]

于是没有定义的合成只有 \(\infty + \infty\)、\(\infty, 0\) 与 \(0, \infty\)。另一方面，映射 \(x \mapsto -x\) 与 \(x \mapsto x^{-1}\) 类似地扩张到 \(\tilde{\mathbf{K}}\)，通过令 \(-\infty = \infty, 0^{-1} = \infty, \infty^{-1} = 0\)。我们还将记 \(x + (-y) = x - y\)。

集 \(\tilde{\mathbf{K}}\) 称为与 \(\mathbf{K}\) 相伴的射影域，它可等同于射影直线 \(\mathbf{P}_1(\mathbf{K})\)（出处同上）。

定义 3. 设 K、L 是两个域。\(\tilde{\mathbf{K}}\) 到 \(\tilde{\mathbf{L}}\) 的（关于加法与乘法的）每个满足 \(f(1) = 1\) 的态射 f 称为 K 的以 L 为值域的位。

换言之，若 \(x\)、\(y\) 是 \(\tilde{\mathbf{K}}\) 的元素，且 \(f(x) + f(y)\)（相应地，\(f(x)f(y)\)）有定义，则 \(x + y\)（相应地，\(xy\)）有定义，且

\begin{enumerate}
\def\labelenumi{(\arabic{enumi})}
\setcounter{enumi}{3}
\tightlist
\item
\end{enumerate}

\[
f (x + y) = f (x) + f (y)\tag{4}
\]

\[
f (x y) = f (x) f (y).\tag{5}
\]

由于 \(\infty +\infty\) 没有定义，\(f(\infty) + f(\infty)\) 也没有定义，这表明

\[
f (\infty) = \infty .\tag{6}
\]

类似地，由于 \(0.\infty\) 没有定义，\(f(0)f(\infty)\) 也没有定义，由 (6) 这蕴涵

\[
f (0) = 0.\tag{7}
\]

另一方面

\[
f (a ^ {- 1}) = f (a) ^ {- 1} \quad \text {   for   all   } a \in \tilde {\mathbf {K}}.\tag{8}
\]

若 \(f(a)f(a^{-1})\) 有定义，则 \(aa^{-1}\) 有定义，从而等于 1；于是 \(f(a)f(a^{-1}) = f(1) = 1\)，这在此情形证明了 (8)。若 \(f(a)f(a^{-1})\) 没有定义，则或者 \(f(a) = 0\) 且 \(f(a^{-1}) = \infty\)，或者 \(f(a) = \infty\) 且 \(f(a^{-1}) = 0\)，而 (8) 仍成立。

类似地可以证明

\[
f (- a) = - f (a) \quad \text {   for   all   } a \in \tilde {\mathbf {K}}.\tag{9}
\]

由公式 (8) 与 (9) 得出，\(f\) 对合成律 \((x, y) \mapsto x - y\) 与 \((x, y) \mapsto xy^{-1}\) 也是态射。

对 \(x \in \tilde{\mathbf{K}}\)，称 \(f\) 在 \(x\) 处有限，如果 \(f(x) \neq \infty\)；由 (6) 这蕴涵 \(x \in \mathbf{K}\)。

若 \(f \colon \tilde{\mathbf{K}} \to \tilde{\mathbf{L}}\) 是位，则 \(f(\tilde{\mathbf{K}})\) 是 \(\tilde{\mathbf{L}}\) 的一个子集，它对合成律 \((x, y) \mapsto x + y\)、\((x, y) \mapsto x - y\)、\((x, y) \mapsto xy\) 与 \((x, y) \mapsto xy^{-1}\) 稳定，且包含 1。若 \(\mathbf{E}\) 是 \(f(\tilde{\mathbf{K}})\) 的有限元素所成之集，则 \(\mathbf{E}\) 是 \(\mathbf{L}\) 的子域，且 \(f(\tilde{\mathbf{K}}) = \tilde{\mathbf{E}}\)。习惯上称 \(\mathbf{E}\) 为 \(f\) 的值域。

两个位的合成映射是位。

设 F 是从域 K 到域 L 的一个子域上的同构；把 f 扩张到 \(\tilde{K}\)，令 \(f(\infty) = \infty\)。如此得到 K 的一个以 L 为值域的位，称之为平凡位，它常与同构 f 等同。

